% !TEX program = lualatex
\documentclass[aspectratio=43,professionalfonts,handout]{beamer}
\usepackage{beamer_template}
\usepackage{enumerate}

\newcommand{\LectureNum}{3}
\newcommand{\LectureTitle}{微分法則と積分法則}
\newcommand{\CourseName}{応用数学}
\renewcommand{\TermName}{前期}

\begin{document}

\begin{frame}{第3回　微分法則と積分法則}
  \begin{block}{本回の内容}
  微分法則と積分法則
  \end{block}
  \vfill\hfill{\scriptsize \textcolor{gray}{第2章 ラプラス変換　§1}}
\end{frame}

\begin{frame}{微分法則と積分法則}
  \begin{block}{}
  {\large 微分法則と積分法則}
  \end{block}
  \vfill\hfill{\scriptsize \textcolor{gray}{p.\,52--57}}
\end{frame}

\begin{frame}{定理：原関数の微分法則}
  \begin{block}{}
  \begin{align*}
    L[f'(t)]    &= sF(s) - f(0) \\
    L[f''(t)]   &= s^{2}F(s) - sf(0) - f'(0) \\
    L[f^{(n)}(t)] &= s^n F(s) - s^{n-1}f(0) - s^{n-2}f'(0) - \ldots - f^{(n-1)}(0)
  \end{align*}
  \end{block}
  \vfill\hfill{\scriptsize \textcolor{gray}{p.\,52--57}}
\end{frame}

\begin{frame}{定理：像関数の高次微分}
  \begin{block}{}
  \[
    L[t^n f(t)] = (-1)^n F^{(n)}(s)
  \]
  \end{block}
  \vfill\hfill{\scriptsize \textcolor{gray}{p.\,52--57}}
\end{frame}

\begin{frame}{定理：原関数の積分法則}
  \begin{block}{}
  \[
    L[\int_{0}^t f(\tau)d\tau] = F(s)/s
  \]
  \end{block}
  \vfill\hfill{\scriptsize \textcolor{gray}{p.\,52--57}}
\end{frame}

\begin{frame}{定理：像関数の積分}
  \begin{block}{}
  \[
    L[f(t)/t] = \int_{s}^\infty F(u)du
  \]
  \end{block}
  \vfill\hfill{\scriptsize \textcolor{gray}{p.\,52--57}}
\end{frame}

\begin{frame}{例題9}
  \begin{exampleblock}{問題}
    微分法則を用いて次の問いに答えよ\\[2pt]
  \end{exampleblock}
\end{frame}

\begin{frame}{例題9　解答}
  \begin{exampleblock}{解答}
  \end{exampleblock}
\end{frame}

\begin{frame}{例題10}
  \begin{exampleblock}{問題}
    $L[t^n]$ を求めよ（ n は正の整数）\\[2pt]
  \end{exampleblock}
\end{frame}

\begin{frame}{例題10　解答}
  \begin{exampleblock}{解答}
  \begin{enumerate}[1.]
    \item[] $f(t) = t^n$ とおくと
    \item[] $f^{(n)}(t) = n!, f(0) = f'(0) = \ldots = f^{(n-1)}(0) = 0$
    \item[] 微分法則を n 回適用 $: s^n F(s) - 0 = L[n!] = n!/s$
  \end{enumerate}
  \end{exampleblock}
  \begin{alertblock}{答}
    $L[t^n] = n!/s^{n+1}  (s > 0)$
  \end{alertblock}
\end{frame}

\begin{frame}{例題11}
  \begin{exampleblock}{問題}
    $L[\sin(\alpha t)/t]$ を求めよ（ $\alpha$ は正の定数）\\[2pt]
  \end{exampleblock}
\end{frame}

\begin{frame}{例題11　解答}
  \begin{exampleblock}{解答}
  \begin{enumerate}[1.]
    \item[] 像関数の積分 $: L[f(t)/t] = \int_{s}^\infty F(u)du$ を利用
    \item[] $L[\sin(\alpha t)] = \alpha/(u^{2}+\alpha^{2})$ なので
    \item[] $L[\sin(\alpha t)/t] = \int_{s}^\infty \alpha/(u^{2}+\alpha^{2}) du = [\arctan(u/\alpha)]_{s}^\infty$
    \item[] $= \pi/2 - \arctan(s/\alpha)$
  \end{enumerate}
  \end{exampleblock}
  \begin{alertblock}{答}
    $L[\sin(\alpha t)/t] = \pi/2 - \arctan(s/\alpha) = \arctan(\alpha/s)  (s > 0)$
  \end{alertblock}
\end{frame}

\begin{frame}{演習問題（p.\,52--57）}
  \begin{exampleblock}{自分で解いてみよう}
  \begin{description}
    \item[問13] $L[t \cos(\alpha t)]$ を求めよ
    \item[問14] $f(t)$ が $f'(t)+2f(t)=f(0)$ を満たすとき、 $L[f(t)]$ を求めよ
    \item[問15] $L[t^n]$ を求めよ（ n は正の整数）（像関数の高次微分を使用）
    \item[問16] $L[(e^{2t} - e^t)/t]$ を求めよ
  \end{description}
  \end{exampleblock}
\end{frame}

\begin{frame}{解説（試験前配布）}
  \vfill
  \begin{center}
  {\LARGE\bfseries\color{structure.fg} 解説（試験前配布）}\\[1em]
  {\normalsize\textcolor{gray}{（試験前に Moodle で配布）}}
  \end{center}
  \vfill
\end{frame}

\begin{frame}{問13　解答}
  \begin{block}{}
  \begin{itemize}\setlength{\itemsep}{2pt}
    \item 像関数の高次微分 $: L[t^n f(t)] = (-1)^n F^{(n)}(s)$ を $n=1$ で使用
    \item $F(s) = L[\cos(\alpha t)] = s/(s^{2}+\alpha^{2})$
    \item $F'(s) = d/ds [s/(s^{2}+\alpha^{2})] = (s^{2}+\alpha^{2}-2s^{2})/(s^{2}+\alpha^{2})^{2} = (\alpha^{2}-s^{2})/(s^{2}+\alpha^{2})^{2}$
    \item $L[t \cos(\alpha t)] = (-1)^{1}\cdot F'(s) = -F'(s)$
  \end{itemize}
  \end{block}
  \begin{exampleblock}{答え}
    \[
    L[t \cos(\alpha t)] = (s^{2}-\alpha^{2})/(s^{2}+\alpha^{2})^{2}
  \]
  \end{exampleblock}
\end{frame}

\begin{frame}{問14　解答}
  \begin{block}{}
  \begin{itemize}\setlength{\itemsep}{2pt}
    \item 両辺をラプラス変換 $: (sF(s)-f(0)) + 2F(s) = f(0)/s$
    \item $(s+2)F(s) = f(0) + f(0)/s = f(0)(s+1)/s$
    \item $F(s) = f(0)(s+1)/(s(s+2))$
    \item 部分分数 $: (s+1)/(s(s+2)) = A/s + B/(s+2)$
    \item $A = 1/2, B = 1/2$
  \end{itemize}
  \end{block}
  \begin{exampleblock}{答え}
    \[
    F(s) = (f(0)/2)(1/s + 1/(s+2))   \to   f(t) = (f(0)/2)(1 + e^{-2t})
  \]
  \end{exampleblock}
  {\footnotesize \textcolor{gray}{\textit{問題文は $f'(t)+2f(t)-f(0)=0 \to f'(t)+2f(t)=f(0)$ と同値}}}
\end{frame}

\begin{frame}{問15　解答}
  \begin{block}{}
  \begin{itemize}\setlength{\itemsep}{2pt}
    \item $f(t) = 1$ とおくと $F(s) = L[1] = 1/s$
    \item $F^{(n)}(s) = d^n/ds^n [1/s] = (-1)^n n!/s^{n+1}$
    \item $L[t^n\cdot1] = (-1)^n F^{(n)}(s) = (-1)^n\cdot(-1)^n n!/s^{n+1}$
  \end{itemize}
  \end{block}
  \begin{exampleblock}{答え}
    \[
    L[t^n] = n!/s^{n+1}  (s > 0)
  \]
  \end{exampleblock}
\end{frame}

\begin{frame}{問16　解答}
  \begin{block}{}
  \begin{itemize}\setlength{\itemsep}{2pt}
    \item 像関数の積分 $: L[f(t)/t] = \int_{s}^\infty F(u)du$ を利用
    \item $f(t) = e^{2t} - e^t$ とおくと $F(u) = L[e^{2t}-e^t]|_{s=u} = 1/(u-2) - 1/(u-1)$
    \item $L[(e^{2t}-e^t)/t] = \int_{s}^\infty (1/(u-2) - 1/(u-1)) du$
    \item $= [\ln(u-2) - \ln(u-1)]_{s}^\infty$
    \item $= \lim_{u \to \infty} \ln((u-2)/(u-1)) - \ln((s-2)/(s-1))$
    \item $= 0 - \ln((s-2)/(s-1)) = \ln((s-1)/(s-2))$
  \end{itemize}
  \end{block}
  \begin{exampleblock}{答え}
    \[
    L[(e^{2t}-e^t)/t] = \ln((s-1)/(s-2))  (s > 2)
  \]
  \end{exampleblock}
  {\footnotesize \textcolor{gray}{\textit{p.56 スキャンで確認（元 YAML の $e^{-\tau}/\tau$ は誤記）}}}
\end{frame}

\end{document}
